\scanpage{10}
\mspage{$6'$}

En plus de ces $5$ variantes $\mathcal X_0,\mathcal X_1,\mathcal X'_1,\mathcal X_2,\mathcal X_3$ il y a \inserted{bien entendu} les revêtements de ces cartes (ramifiés pour les trois premières) principaux de groupe $D_{2n}$, associés aux $\operatorname{Pol}(D)$. Ce sont les revêtements $\widetilde{\mathcal X}_i$ qui sont les déploiements ultimes.

Les surfaces les plus utiles me semblent ici $\mathcal X'_1,\mathcal X_2,\mathcal X_3$ et leurs revêtements $\widetilde{\mathcal X}'_1,\widetilde{\mathcal X}_2,\widetilde{\mathcal X}_3$. Peut-être peut-on se passer des $\mathcal X_2,\mathcal X_3$ et leurs tildés, au profit des $\overset{\deleted{\unsure{\sim}}}{\mathcal X}'_1$ et $\widetilde{\mathcal X}'_1$, d'autant que $\mathcal X'_1$ s'envoie bien dans $X$ comme prévu --- je \inserted{ne} suis pas tout à fait \unsure{convaincu} des \unsure{collisions|collages}, \deleted{\unsure{ayant quelque}}\inserted{en usage} \unsure{qd}, comment ça se passe au voisinage du bord d'une face supercritique.

\margin{9}

\begin{tabular}{@{}l@{}}
\texttt{********************************************************************************}\\
\texttt{*JESNEWS\ DU\ 02/08/1982\ \unsure{C63|C68}\ CADARACHE\ \ \ \ \ \ \ \ \ \ \ \ \ \ \ \ \ \ \ \ \ \ \ \ \ \ \ \ \ \ \ \ \ \ \ \ \ \ \ \ \ \ \ \ \ *}\\
\texttt{********************************************************************************}\\
\texttt{*\ \ \ \ \ \ DEPUIS\ LE\ 1/7/82\ LES\ ALIAS\ DES\ COMMANDES\ LIBRARIAN\ \ \ \ \ \ \ \ \ \ \ \ \ \ \ \ \ \ \ \ \ \ *}\\
\texttt{*\ \ \ \ \ \ \ \ \ \ \ \ \ LIBGET\ (LG)\ ET\ LIBSAVE\ (LS)\ SONT\ SUPPRIMES,\ \ \ \ \ \ \ \ \ \ \ \ \ \ \ \ \ \ \ \ \ *}\\
\texttt{*\ \ \ \ \ \ \ \ \ \ \ \ \ POUR\ EVITER\ UNE\ POSSIBILITE\ DE\ CONFLIT\ AVEC\ \ \ \ \ \ \ \ \ \ \ \ \ \ \ \ \ \ \ \ *}\\
\texttt{*\ \ \ \ \ \ \ \ \ \ \ \ \ LE\ MEME\ ALIAS\ DES\ COMMANDES\ STANDARD\ IBM.\ \ \ \ \ \ \ \ \ \ \ \ \ \ \ \ \ \ \ \ \ *}\\
\texttt{*\ \ \ \ \ \ ********\ \ \ \ \ \ \ \ \ \ \ \ \ \ \ \ \ \ \ \ \ \ \ \ \ \ \ \ \ \ \ \ \ \ \ \ \ \ \ \ \ \ \ \ \ \ \ \ \ \ \ \ \ \ \ \ \ \ \ \ \ *}\\
\texttt{*\ ATTENTION\ \ \ \ \ \ \ \ \ \ \ \ \ \ \ \ \ \ \ \ \ \ \ \ \ \ \ \ \ \ \ \ \ \ \ \ \ \ \ \ \ \ \ \ \ \ \ \ \ \ \ \ \ \ \ \ \ \ \ \ \ \ \ \ \ *}\\
\texttt{*\ \ \ MISE\ EN\ SERVICE\ LE\ VENDREDI\ 30/7/82\ A\ 8H00\ DE\ LA\ DEUXIEME\ PHASE\ DE\ G4F\ \ \ \ *}\\
\texttt{*\ \ \ \ \ SUR\ 168\ CADARACHE.\ \ \ \ \ \ \ \ \ \ \ \ \ \ \ \ \ \ \ \ \ \ \ \ \ \ \ \ \ \ \ \ \ \ \ \ \ \ \ \ \ \ \ \ \ \ \ \ \ \ \ \ \ *}\\
\texttt{*\ \ \ \ \ ELLE\ PERMET\ D'UTILISER\ DANS\ LA\ CDE.\ KEEP\ TSO\ ET\ DANS\ LA\ PROCEDURE\ KEEP\ \ *}\\
\texttt{*\ \ \ \ \ BATCH\ UNE\ POSSIBILITE\ D'AFFECTATION\ DE\ TYPE\ DE\ SAUVEGARDE\ AU\ NIVEAU\ DU\ *}\\
\texttt{*\ \ \ \ \ FICHIER.CETTE\ POSSIBILITE\ EST\ VIVEMENT\ CONSEILLEE.\ \ \ \ \ \ \ \ \ \ \ \ \ \ \ \ \ \ \ \ \ \ *}\\
\texttt{*\ \ \ \ \ POUR\ PLUS\ AMPLES\ INFORMATIONS\ VOIR\ NOTE\ TECHNIQUE\ CISI\ MVS140,LA\ NOTE\ \ *}\\
\texttt{*\ \ \ \ \ MVS\ 143\ ETANT\ SUPPRIMEE.\ \ \ \ \ \ \ \ \ \ \ \ \ \ \ \ \ \ \ \ \ \ \ \ \ \ \ \ \ \ \ \ \ \ \ \ \ \ \ \ \ \ \ \ \ \ \ *}\\
\texttt{*\ \ \ \ \ \ \ \ \ \ \ \ \ \ \ \ \ \ \ \ \ \ \ \ \ \ \ \ \ \ \ \ \ \ \ \ \ \ \ \ \ \ \ \ \ \ \ \ \ \ \ \ \ \ \ \ \ \ \ \ \ \ \ \ \ \ \ \ \ \ \ \ \ \ \ \ *}\\
\texttt{*\ EXEMPLES\ D'UTILISATION:\ \ \ \ \ \ \ \ \ \ \ \ \ \ \ \ \ \ \ \ \ \ \ \ \ \ \ \ \ \ \ \ \ \ \ \ \ \ \ \ \ \ \ \ \ \ \ \ \ \ \ \ *}\\
\texttt{*\ \ \ \ \ \ \ \ \ \ \ \ \ \ \ \ \ \ \ \ \ \ \ \ \ \ \ \ \ \ \ \ \ \ \ \ \ \ \ \ \ \ \ \ \ \ \ \ \ \ \ \ \ \ \ \ \ \ \ \ \ \ \ \ \ \ \ \ \ \ \ \ \ \ \ \ *}\\
\texttt{*\ \ BATCH:\ \ DSN=NOM.DU.FICHIER.A.SAUVEGARDER,DELAY=TYPE\ DE\ SAUVEGARDE\ \ \ \ \ \ \ \ \ \ \ *}\\
\texttt{*\ \ \ \ \ \ \ \ \ \ \ \ \ \ \ \ \ \ \ \ \ \ \ \ \ \ \ \ \ \ \ \ \ \ \ \ \ \ \ \ \ \ \ \ \ \ \ \ \ \ \ \ \ \ \ \ \ \ \ \ \ \ \ \ \ \ \ \ \ \ \ \ \ \ \ \ *}\\
\texttt{*\ \ TSO:\ \ \ \ \ \ \ \ \ \ \ \ \ \ \ \ \ \ \ \ \ \ \ \ \ \ \ \ \ \ \ \ \ \ \ \ \ \ \ \ \ \ \ \ \ \ \ \ \ \ \ \ \ \ \ \ \ \ \ \ \ \ \ \ \ \ \ \ \ *}\\
\texttt{*\ \ K\ DA('NOM\ DU\ FICHIER\ A\ SAUVEGARDER')\ DELAY(TYPE\ DE\ SAUVEGARDE)\ \ \ \ \ \ \ \ \ \ \ \ \ *}\\
\texttt{*\ \ \ \ \ \ \ \ \ \ \ \ \ \ \ \ \ \ \ \ \ \ \ \ \ \ \ \ \ \ \ \ \ \ \ \ \ \ \ \ \ \ \ \ \ \ \ \ \ \ \ \ \ \ \ \ \ \ \ \ \ \ \ \ \ \ \ \ \ \ \ \ \ \ \ \ *}\\
\texttt{*\ \ TYPES\ DE\ SAUVEGARDE:\ \ \ \ \ \ \ \ \ \ \ \ \ \ \ \ \ \ \ \ \ \ \ \ \ \ \ \ \ \ \ \ \ \ \ \ \ \ \ \ \ \ \ \ \ \ \ \ \ \ \ \ \ \ *}\\
\texttt{*\ \ \ \ \ \ \ \ \ \ \ \ \ \ \ \ \ \ \ \ \ \ \ \ \ \ \ \ \ \ \ \ \ \ \ \ \ \ \ \ \ \ \ \ \ \ \ \ \ \ \ \ \ \ \ \ \ \ \ \ \ \ \ \ \ \ \ \ \ \ \ \ \ \ \ \ *}\\
\texttt{*\ \ \ \ \ \ W\ \ \ \ \ \ \ SAUVEGARDE\ HEBDO.\ \ \ \ \ \ \ \ \ \ \ \ \ \ \ \ \ \ \ \ \ \ \ \ \ \ \ \ \ \ \ \ \ \ \ \ \ \ \ \ \ \ \ \ \ *}\\
\texttt{*\ \ \ \ \ \ D\ \ \ \ \ \ \ SAUVEGARDE\ JOURNALIERE.\ \ \ \ \ \ \ \ \ \ \ \ \ \ \ \ \ \ \ \ \ \ \ \ \ \ \ \ \ \ \ \ \ \ \ \ \ \ \ \ *}\\
\texttt{*\ \ \ \ \ \ NOS\ \ \ \ \ SAUVEGARDE\ MENSUELLE.\ \ \ \ \ \ \ \ \ \ \ \ \ \ \ \ \ \ \ \ \ \ \ \ \ \ \ \ \ \ \ \ \ \ \ \ \ \ \ \ \ \ *}\\
\texttt{*\ \ \ \ \ \ *\ \ \ \ \ \ \ PARAMETRE\ PAR\ DEFAUT.\ \ \ \ \ \ \ \ \ \ \ \ \ \ \ \ \ \ \ \ \ \ \ \ \ \ \ \ \ \ \ \ \ \ \ \ \ \ \ \ \ *}\\
\texttt{*\ \ \ \ \ \ \ \ \ \ \ \ \ \ \ \ \ \ \ \ \ \ \ \ \ \ \ \ \ \ \ \ \ \ \ \ \ \ \ \ \ \ \ \ \ \ \ \ \ \ \ \ \ \ \ \ \ \ \ \ \ \ \ \ \ \ \ \ \ \ \ \ \ \ \ \ *}\\
\texttt{********************************************************************************}\\
\\
\texttt{*******************************************************************}\\
\texttt{*\ \ \ \ \ \ \ \ \ \ \ \ \ \ \ \ \ \ \ \ \ \ \ \ \ \ \ \ \ \ \ \ \ \ \ \ \ \ \ \ \ \ \ \ \ \ \ \ \ \ \ \ \ \ \ \ \ \ \ \ \ \ \ \ \ *}\\
\texttt{*\ \ \ \ \ \ \ \ \ \ \ \ \ \ \ \ \ \ -\ J\ E\ S\ N\ E\ W\ S\ \ F\ R\ 1\ -\ \ \ \ \ \ LE\ 06/08/82\ \ \ \ \ \ *}\\
\texttt{*\ \ \ \ \ \ \ \ \ \ \ \ \ \ \ \ \ \ \ \ \ \ \ \ \ \ \ \ \ \ \ \ \ \ \ \ \ \ \ \ \ \ \ \ \ \ \ \ \ \ \ \ \ \ \ \ \ \ \ \ \ \ \ \ \ *}\\
\texttt{*-----------------------------------------------------------------*}\\
\texttt{*\ \ \ \ \ \ \ \ \ \ \ \ \ \ \ \ \ \ \ \ \ \ \ \ \ \ \ \ \ \ \ \ \ \ \ \ \ \ \ \ \ \ \ \ \ \ \ \ \ \ \ \ \ \ \ \ \ \ \ \ \ \ \ \ \ *}\\
\texttt{*\ \ DEPUIS\ LE\ LUNDI\ 12\ JUILLET\ 82,\ LES\ CLASSES\ DE\ TRAVAUX\ D\ ET\ E\ \ \ \ \ \ *}\\
\texttt{*\ \ SERONT\ ETENDUES\ POUR\ UNE\ TAILLE\ MEMOIRE\ SUPERIEURE\ A\ 1024\ K\ \ \ \ \ \ \ *}\\
\texttt{*\ \ SE\ REPORTER\ A\ LA\ NOTE\ TECHNIQUE\ MVS141.\ \ \ \ \ \ \ \ \ \ \ \ \ \ \ \ \ \ \ \ \ \ \ \ \ \ *}\\
\texttt{*\ \ \ \ \ \ \ \ \ \ \ \ \ \ \ \ \ \ \ \ \ \ \ \ \ \ \ \ \ \ \ \ \ \ \ \ \ \ \ \ \ \ \ \ \ \ \ \ \ \ \ \ \ \ \ \ \ \ \ \ \ \ \ \ \ *}\\
\texttt{*-----------------------------------------------------------------*}\\
\texttt{*\ \ A\ PARTIR\ DU\ LUNDI\ 9\ AOUT\ 1982\ ,\ IL\ SERA\ POSSIBLE\ D'IMPRIMER\ \ \ \ \ \ \ *}\\
\texttt{*\ \ SUR\ XEROX\ 9700\ (DESTINATION\ SAC)\ DES\ FICHIERS\ UTILISANT\ LE\ \ \ \ \ \ \ \ *}\\
\texttt{*\ \ COUPLE\ DE\ JEUX\ DE\ CARACTERES\ :\ (GS15,GT15)\ \ \ \ \ \ \ \ \ \ \ \ \ \ \ \ \ \ \ \ \ \ \ \ *}\\
\texttt{*******************************************************************}\\
\end{tabular}
